% BHU PYQ -- Professional Exam Template
\documentclass[11pt,a4paper]{article}

\usepackage[a4paper,top=2.5cm,bottom=2.5cm,left=2.8cm,right=2.8cm,headheight=14pt]{geometry}
\usepackage{amsmath,amssymb}
\usepackage[T1]{fontenc}
\usepackage[utf8]{inputenc}
\usepackage{lmodern}
\usepackage{microtype}
\usepackage{fancyhdr}
\usepackage{enumitem}
\usepackage{listings}
\usepackage{hyperref}
\usepackage{lastpage}
\usepackage{parskip}

\hypersetup{colorlinks=true,urlcolor=black,linkcolor=black,
  pdftitle={CS-108: Data Structures Algorithms -- BHU PYQ}}

\pagestyle{fancy}
\fancyhf{}
\renewcommand{\headrulewidth}{0.4pt}
\renewcommand{\footrulewidth}{0.4pt}
\fancyhead[L]{\small\textit{Banaras Hindu University}}
\fancyhead[R]{\small\textit{CS-108: Data Structures Algorithms}}
\fancyfoot[C]{\small Page \thepage\ of \pageref{LastPage}}

\lstset{
  basicstyle=\small\ttfamily,
  frame=single,
  breaklines=true,
  columns=flexible,
  keepspaces=true
}

\newlist{parts}{enumerate}{2}
\setlist[parts,1]{label=(\alph*),leftmargin=*,itemsep=4pt,topsep=3pt}
\setlist[parts,2]{label=(\roman*),leftmargin=*,itemsep=2pt,topsep=2pt}
\newcommand{\pts}[1]{\hfill\textit{[#1]}}

\begin{document}

\begin{center}
  {\large\bfseries BANARAS HINDU UNIVERSITY}\\[2pt]
  {\normalsize B.Sc. (Hons.) Semester V Examination, 2022-23}\\[6pt]
  \rule{0.8\linewidth}{0.5pt}\\[6pt]
  {\large\bfseries Paper: CS-108 --- Data Structures Algorithms}\\[4pt]
  {\normalsize Time: Three hours \hfill Maximum Marks: 70}
\end{center}

\rule{\linewidth}{0.4pt}

\smallskip
\noindent\textit{Instructions: Attempt any five questions. All questions carry equal marks.}

\bigskip

\medskip\noindent\textbf{Question 1.}
\begin{parts}
  \item Explain primitive and non-primitive data structures with suitable examples.
  \item Distinguish between arrays and linked lists with suitable illustrations.
\end{parts}

\medskip\noindent\textbf{Question 2.}
\begin{parts}
  \item What is a Binary Search Tree (BST)? How is a BST different from a B-Tree? Explain with a suitable example.
  \item Discuss tree traversal techniques with suitable examples.
\end{parts}

\medskip\noindent\textbf{Question 3.}
\begin{parts}
  \item Define Quick Sort. Sort the following array using the Quick Sort algorithm:
  \[ 12, 98, 1, 55, 3, 77, 0, 10 \]
  \item Discuss the time complexity of Quick Sort in the best, average, and worst cases.
\end{parts}

\medskip\noindent\textbf{Question 4.}
\begin{parts}
  \item Explain Greedy and Dynamic programming approaches.
  \item Choose a problem of your choice and solve it using both approaches.
\end{parts}

\medskip\noindent\textbf{Question 5.}
\begin{parts}
  \item Define asymptotic notations ($\mathcal{O}, \Omega, \Theta$) with suitable examples.
  \item Find the upper, lower, and tight bounds of the running time of the quadratic function:
  \[ f(n) = 99n^2 + 28n + 4 \]
\end{parts}

\medskip\noindent\textbf{Question 6.}
\begin{parts}
  \item Define backtracking. Discuss it through a suitable example.
  \item Define Graph as a data structure. Explain the adjacency matrix and adjacency list representations with suitable examples.
\end{parts}

\medskip\noindent\textbf{Question 7.}
\begin{parts}
  \item Define the shortest path problem with suitable illustrations.
  \item Define Hashing and Collision. What are the methods used to handle collisions? Explain with an example.
\end{parts}

\vfill
\rule{\linewidth}{0.4pt}
\begin{center}\small
\href{https://bhu-exam.vercel.app/}{Click here to visit our website}\\[4pt]\href{https://me-aryan.vercel.app/}{Click here to visit our contributor}\\[4pt]\href{https://quashan-me.vercel.app/}{Click here to visit our contributor}
\end{center}

\end{document}
